% Pacotes e ajustes de formatação ao Guia UECE 2026.
% Páginas citadas são do PDF do guia (página impressa = PDF − 2).

\usepackage[T1]{fontenc}
\usepackage[utf8]{inputenc}
\usepackage{newtxtext,newtxmath}   % Times (guia p. 76: Arial ou Times 12)
\usepackage{indentfirst}
\usepackage{xcolor}
\usepackage{graphicx}
\usepackage{booktabs}              % tabelas com fios só no cabeçalho e fechamento (p. 85-86, IBGE)
\usepackage{longtable}
\usepackage{multirow}
\usepackage{listings}              % SQL e YAML nos exemplos e apêndices
\usepackage{tikz}
\usetikzlibrary{positioning,arrows.meta,fit,calc}
\usepackage{array}
\usepackage{tabularx}
\usepackage{microtype}
\usepackage[brazilian,hyperpageref]{backref}
\usepackage[alf,abnt-emphasize=bf]{abntex2cite}  % autor-data (p. 103, 112)

% --- Marcação de pendências -------------------------------------------------
\newcommand{\pendente}[1]{\textcolor{red}{[PENDENTE: #1]}}
% Roteiro de rascunho: \begin{esboco} ... \end{esboco} (cada um em linha
% própria). Conteúdo lido verbatim e descartado — não aparece no PDF.
% Para revisar o roteiro no PDF, troque \excludecomment por \includecomment
% (aí o conteúdo precisa ser LaTeX válido).
\usepackage{comment}
\excludecomment{esboco}

% --- Nomes de arquivo, tabela e coluna (aceita "_" sem escape) ------------
\newcommand{\arq}[1]{\texttt{\detokenize{#1}}}

% --- Margens (p. 76): sup. e esq. 3 cm; dir. e inf. 2 cm ---------------------
\setlrmarginsandblock{3cm}{2cm}{*}
\setulmarginsandblock{3cm}{2cm}{*}
\checkandfixthelayout

% --- Parágrafo e espaçamento (p. 76-79) --------------------------------------
\setlength{\parindent}{2cm}
\setlength{\parskip}{0pt}
\OnehalfSpacing

% --- Hierarquia de títulos (p. 61, 79-80), todos em corpo 12 ----------------
% primária: MAIÚSCULAS + negrito | secundária: negrito | terciária: normal
% quaternária: negrito + itálico | quinária: itálico
\renewcommand{\ABNTEXchapterfont}{\bfseries}
\renewcommand{\ABNTEXchapterfontsize}{\normalsize}
\renewcommand{\ABNTEXsectionfont}{\bfseries}
\renewcommand{\ABNTEXsectionfontsize}{\normalsize}
\renewcommand{\ABNTEXsectionupperifneeded}[1]{#1}
\renewcommand{\ABNTEXsubsectionfont}{\normalfont}
\renewcommand{\ABNTEXsubsectionfontsize}{\normalsize}
\renewcommand{\ABNTEXsubsectionupperifneeded}[1]{#1}
\renewcommand{\ABNTEXsubsubsectionfont}{\bfseries\itshape}
\renewcommand{\ABNTEXsubsubsectionfontsize}{\normalsize}
\renewcommand{\ABNTEXsubsubsectionupperifneeded}[1]{#1}
\renewcommand{\ABNTEXsubsubsubsectionfont}{\itshape}
\renewcommand{\ABNTEXsubsubsubsectionfontsize}{\normalsize}

% --- Paginação (p. 79-80): arábico, canto sup. direito, corpo 10 ------------
\makeoddhead{abntheadings}{}{}{\footnotesize\thepage}

% --- Legendas (p. 85): título acima, centralizado, negrito; fonte abaixo -----
\captionnamefont{\bfseries\normalsize}
\captiontitlefont{\bfseries\normalsize}
\captiondelim{ -- }
% \fonte{...} do abnTeX2 já posiciona a fonte abaixo, em corpo menor.

% --- Quadros (dados descritivos, fechados; p. 50) ---------------------------
\newcommand{\listofquadrosname}{Lista de quadros}
\newcommand{\quadroname}{Quadro}
\newfloat[chapter]{quadro}{loq}{\quadroname}
\newlistof{listofquadros}{loq}{\listofquadrosname}
\newlistentry{quadro}{loq}{0}
\renewcommand{\cftquadroname}{\quadroname\space}
\renewcommand*{\cftquadroaftersnum}{\hfill--\hfill}

% --- Nomes usados por \autoref (guia p. 61: não usar "capítulo") -----------
\addto\extrasbrazil{%
  \renewcommand{\chapterautorefname}{Seção}%
  \renewcommand{\sectionautorefname}{Seção}%
  \renewcommand{\subsectionautorefname}{Seção}%
  \renewcommand{\subsubsectionautorefname}{Seção}%
}
\providecommand{\quadroautorefname}{Quadro}
\renewcommand{\lstlistingname}{Código}

% --- Código ------------------------------------------------------------------
\lstset{
  basicstyle=\ttfamily\footnotesize,
  breaklines=true,
  columns=fullflexible,
  frame=tb,
  showstringspaces=false,
  inputencoding=utf8,
  extendedchars=true,
  literate={á}{{\'a}}1 {à}{{\`a}}1 {ã}{{\~a}}1 {â}{{\^a}}1 {é}{{\'e}}1
           {ê}{{\^e}}1 {í}{{\'i}}1 {ó}{{\'o}}1 {õ}{{\~o}}1 {ô}{{\^o}}1
           {ú}{{\'u}}1 {ç}{{\c c}}1 {Á}{{\'A}}1 {É}{{\'E}}1 {Ç}{{\c C}}1
}

% --- PDF ----------------------------------------------------------------------
\makeatletter
\hypersetup{
  pdftitle={\@title},
  pdfauthor={\@author},
  pdfsubject={Trabalho de Conclusão de Curso},
  pdfkeywords={text-to-SQL}{camada semântica}{modelos de linguagem}{agentes},
  colorlinks=true,
  linkcolor=black,
  citecolor=black,
  urlcolor=black,
  bookmarksdepth=4
}
\makeatother

\makeindex
